\documentclass[11pt,letterpaper]{article}

\usepackage[T1]{fontenc}
\usepackage{amsmath,amssymb}
\usepackage{tikz}
\usepackage{pgfplots}
\pgfplotsset{compat=1.18}
\usepackage{algorithm}
\usepackage{algpseudocode}
\usepackage{booktabs}
\usepackage{enumitem}
\usepackage[hidelinks]{hyperref}

\newcommand{\N}{\mathbb{N}}
\newcommand{\Z}{\mathbb{Z}}
\newcommand{\Q}{\mathbb{Q}}
\newcommand{\R}{\mathbb{R}}
\newcommand{\C}{\mathbb{C}}
\newcommand{\Map}[3]{#1 \colon #2 \to #3}
\newcommand{\Mod}[1]{#1\text{-}\mathrm{Mod}}
\newcommand{\GroupHom}{\mathrm{Hom}_{\mathbf{Grp}}}
\newcommand{\RingHom}{\mathrm{Hom}_{\mathbf{Ring}}}

\newtheorem{theorem}{Theorem}
\newtheorem{lemma}[theorem]{Lemma}
\newtheorem{corollary}[theorem]{Corollary}
\newtheorem{definition}[theorem]{Definition}
\newtheorem{remark}[theorem]{Remark}
\newenvironment{proof}{\noindent\textit{Proof.}}{\hfill\(\square\)}

\title{Affine Escape from a Prime-Square Residue Layer}
\author{}
\date{\today}

\begin{document}
\maketitle

\begin{abstract}
For every prime \(p\) and positive integers \(Q,C\) not divisible by \(p\), consider the raw map with multiplier \(pQ\), offset \(pC\), and divisor \(p^2\). A residue layer of valuation one is invariant and division-free. The complement of its basin has natural density zero. The proof counts finite valuation words and controls the two truncation errors before passing to the density limit. In particular, density-one escape can hold when the common divisor of the multiplier and divisor divides the offset.
\end{abstract}

\section{The raw map and the invariant layer}

Fix a prime \(p\), positive integers \(Q,C\), and assume \(p\nmid Q C\).
All seeds in this article are positive integers. Define
\[
T(n)=
\begin{cases}
n/p^2,&p^2\mid n,\\
pQn+pC,&p^2\nmid n.
\end{cases}
\]
Every orbit remains positive. The division branch has precedence.
Write \(v_p(n)\) for the exponent of \(p\) in a positive integer \(n\).

\begin{lemma}[Invariant layer]
\label{lem:layer}
If \(v_p(n)=1\), then every raw step is affine and every later state has valuation one. The orbit tends to infinity along the full sequence.
\end{lemma}

\begin{proof}
Since \(p^2\nmid n\), the next state is \(p(Qn+C)\).
As \(p\mid n\) and \(p\nmid C\), the factor \(Qn+C\) is not divisible by \(p\).
The new valuation is exactly one. This proves invariance by induction.
In this layer the exact trajectory is
\[
T^j(n)=(pQ)^j n+pC\sum_{i=0}^{j-1}(pQ)^i.
\]
Here \(pQ\ge2\), so these positive states tend to infinity. No compressed-clock argument is needed after layer entry.
\end{proof}

Every seed has a unique form \(n=p^{2k}u\), with \(k\ge0\) and \(p^2\nmid u\).
Indeed repeated division by \(p^2\) strictly decreases a positive integer.
The first \(k\) raw steps recover \(u\). If \(p\mid u\), Lemma~\ref{lem:layer} applies immediately.
It remains to study units \(p\nmid u\).

\begin{lemma}[Unit return]
\label{lem:return}
Let \(p\nmid u\) and put \(v=v_p(Qu+C)\).
If \(v\) is even, including \(v=0\), the next affine step followed by all forced divisions reaches the invariant layer.
If \(v\) is odd, write \(v=2a+1\). The same block returns to the unit
\[
G(u)=\frac{Qu+C}{p^v}
\]
in exactly \(a+2\) raw steps.
\end{lemma}

\begin{proof}
The affine image \(p(Qu+C)\) has valuation \(v+1\).
Each legal division by \(p^2\) removes exactly two from that valuation.
An odd value of \(v+1\) leaves valuation one; an even value leaves valuation zero.
For \(v=2a+1\), there are \(a+1\) forced divisions after the single affine step.
Their quotient is
\[
\frac{p(Qu+C)}{p^{2a+2}}=\frac{Qu+C}{p^{2a+1}},
\]
a unit. Positivity makes \(Qu+C\) nonzero, so \(v\) is finite.
\end{proof}

\section{The exact finite-word law}

The relative density of a subset of units means its density among positive integers not divisible by \(p\).
For a prescribed word \(\boldsymbol v=(v_1,\ldots,v_h)\) of positive valuations, put \(s=v_1+\cdots+v_h\).
The successive unit quotients are required to have exactly these valuations; this word condition is not merely divisibility by \(p^s\).

\begin{lemma}[Valuation cylinders]
\label{lem:cylinder}
The units having a prescribed finite valuation word \(\boldsymbol v\) form exactly \(p-1\) residue classes modulo \(p^{s+1}\).
Their ordinary density is \((p-1)/p^{s+1}\), and their relative unit density is \(p^{-s}\).
\end{lemma}

\begin{proof}
Work first in \(\Z_p\), defined as the inverse limit of \(\Z/p^j\Z\).
Multiplication by \(Q\) is invertible at every finite level because \(p\nmid Q\), and these inverses are compatible.
For \(v\ge1\), the inverse branch
\[
I_v(y)=\frac{p^v y-C}{Q}
\]
sends every \(p\)-adic unit to a unit: its residue modulo \(p\) is \(-C/Q\ne0\).
It satisfies \(QI_v(y)+C=p^v y\), so the valuation is exactly \(v\).
Conversely any unit with that valuation is obtained from this inverse branch.

The composition of \(h\) inverse branches has the form
\[
I_{v_1}\circ\cdots\circ I_{v_h}(y)
=\frac{p^s y-B}{Q^h}
\]
for an integer \(B\); this follows by induction, multiplying the previous numerator by \(Q\) and adding the next translated term.
As \(y\) varies over the \(p-1\) nonzero residues modulo \(p\), the displayed expression gives \(p-1\) distinct source residues modulo \(p^{s+1}\).
Varying the higher digits of \(y\) fills each of those classes because \(Q^h\) is a unit.

Conversely every source in these classes has the specified valuations.
In fact the inverse construction recovers every intermediate unit and the identities \(Q u_{i-1}+C=p^{v_i}u_i\).
For positive integer sources those identities are exact integral divisions, so the \(p\)-adic cylinder condition recovers the arithmetic unit-quotient itinerary. For words consisting entirely of odd valuations, Lemma~\ref{lem:return} identifies those quotients with the actual finite raw unit-return itinerary.
The density of a residue class modulo \(p^{s+1}\) is \(p^{-s-1}\).
Divide the ordinary density by the density \((p-1)/p\) of the units to obtain \(p^{-s}\).
\end{proof}

Thus the positive valuations have weights \(p^{-v}\).
For \(p>2\) they do not exhaust the unit law: valuation zero has the remaining weight \((p-2)/(p-1)\) and immediately enters the invariant layer.
The total weight of the allowed odd valuations is
\[
\lambda=\sum_{a\ge0}p^{-(2a+1)}=\frac{p}{p^2-1}<1.
\]

\begin{lemma}[Finite truncation]
\label{lem:truncation}
Let \(U_h\) be the set of units whose first \(h\) blocks all return to units.
Its upper relative density is at most \(\lambda^h\).
\end{lemma}

\begin{proof}
Fix \(h\) and an integer \(V\ge1\).
Words with all valuations odd and at most \(V\) form a finite disjoint union of the cylinders in Lemma~\ref{lem:cylinder}.
Its relative density is
\[
\left(\sum_{\substack{1\le v\le V\\v\ {\rm odd}}}p^{-v}\right)^h
\le\lambda^h.
\]
If a seed of \(U_h\) is omitted, let \(j<h\) be the first block whose valuation exceeds \(V\).
The preceding \(j\) valuations are odd and at most \(V\), so there are only finitely many preceding words.
For each such word the next condition is \(Qu_j+C\equiv0\pmod{p^{V+1}}\).
Among units its relative density is \(p^{-V}/(p-1)\): it is one unit class modulo \(p^{V+1}\).
The inverse branches multiply this density by the preceding word's weight.
Summing those weights gives at most \(\lambda^j\le1\).
The omitted set therefore lies in a finite union with upper relative density at most
\[
\sum_{j<h}\frac{p^{-V}}{p-1}\le h p^{-V}.
\]
Consequently \(\overline d_{\rm unit}(U_h)\le\lambda^h+h p^{-V}\).
Let \(V\) tend to infinity with \(h\) fixed.
This is a finite-cylinder counting argument; measure zero in \(\Z_p\) alone would not give an ordinary-density statement for integers.
\end{proof}

\section{Density-one escape}

Let \(A_h\) be the positive seeds which, after stripping the initial whole \(p^2\)-power, begin with a unit and survive \(h\) unit returns.

\begin{theorem}[Prime-square leakage]
\label{thm:escape}
For every \(h\ge0\),
\[
\overline d(A_h)\le \frac{p}{p+1}\lambda^h,
\qquad \lambda=\frac{p}{p^2-1}.
\]
The seeds which never enter the valuation-one layer have natural density zero.
The positive seeds whose raw orbits tend to infinity have natural density one.
\end{theorem}

\begin{proof}
Fix an initial-factor cutoff \(K\).
For \(0\le k<K\), multiplication of a subset of integers by \(p^{2k}\) divides its upper ordinary density by \(p^{2k}\):
its counting function at \(X\) is the original counting function at \(\lfloor X/p^{2k}\rfloor\).
Lemma~\ref{lem:truncation}, and the unit density \((p-1)/p\), bound the finite union by
\[
\frac{p-1}{p}\lambda^h\sum_{k=0}^{K-1}p^{-2k}.
\]
Every remaining seed is divisible by \(p^{2K}\); this tail has density \(p^{-2K}\).
Let \(X\) tend to infinity first, then \(K\) tend to infinity. The geometric series gives
\[
\overline d(A_h)\le
\frac{(p-1)/p}{1-p^{-2}}\lambda^h
=\frac{p}{p+1}\lambda^h.
\]

A positive seed which never reaches the layer must survive every number \(h\) of unit returns.
Initial stripping terminates, and each unit block is finite by Lemma~\ref{lem:return}.
Its exceptional set is therefore contained in \(A_h\) for every \(h\).
Since \(\lambda<1\), letting \(h\) tend to infinity proves upper density zero, and hence natural density zero.
Every other seed reaches the layer at a finite raw time and escapes by Lemma~\ref{lem:layer}.
The escaping set contains this density-one set and has density at most one, so its natural density is one.
\end{proof}

For \((q,c,D)=(6,2,4)\), take \(p=2,Q=3,C=1\).
Here \(\gcd(q,D)=2\mid c\), so the common-divisor miss criterion does not apply.
The new surviving weight is \(2/3\).
The conclusion is not an all-seed escape theorem:
for \((2,2,4)\), which satisfies the same hypotheses, \(1\mapsto4\mapsto1\) is a primitive raw two-cycle.
The theorem leaves the zero-density unit-return exceptional set unclassified.

\section{Formal companions}

The active declarations \texttt{SquareLayerImageMiss}, \texttt{SquareLayerInvariant},
\texttt{SquareLayerTail}, \texttt{SquareLayerEscape}, and \texttt{SquareUnitReturn}
in \texttt{QCD.ResidueLeakage.001/lean/NonCoprime.lean} formalize the invariant layer, its full-clock escape, and prescribed odd-power block recovery with a unit endpoint.
The cylinder counting, truncation, and natural-density theorem above are written proofs; these measure-theoretic statements are not asserted to be declarations in that file.

\end{document}
